% language: swedish
% figure: fig1 0.59 0.31 0.72 0.375
% figure: fig2 0.12 0.425 0.41 0.545
\subsection{Vektorer och skalärer}

Kursbok: Vector Calculus, finns som PDF.
Kommer att finnas dugga och inlämningsuppgifter i LP3. Inget i LP2.

\begin{itemize}
  \item Vektorer -- storlek, riktning
  \item Skalär -- storlek
\end{itemize}

Representation av vektor \notefigure[0.15]{fig1}{}

Koordinatsystem

\begin{minipage}{0.32\linewidth}
\notefigure[0.9]{fig2}{}
\end{minipage}
\begin{minipage}{0.66\linewidth}
Introducera vektorbeteckning $\vb{a} = (a_1, a_2, a_3)$
\[ \Rightarrow a_1 = x_2 - x_1, \quad a_2 = y_2 - y_1, \quad a_3 = z_2 - z_1 \]
Högersystem
\end{minipage}

Koordinatsystem beskrivs med enhetsvektorer $\vb{e}_1, \vb{e}_2, \vb{e}_3$ \ $(\hat{x}, \hat{y}, \hat{z})$\\
Basvektorer (ortonormerade)\\
$|\vb{e}_1| = |\vb{e}_2| = |\vb{e}_3| = 1$\\
Kan nu skriva $\vb{a} = a_1\vb{e}_1 + a_2\vb{e}_2 + a_3\vb{e}_3$
\[ |\vb{a}| = \sqrt{a_1^2 + a_2^2 + a_3^2} \]
Ortsvektorer: En vektor som pekar från en plats i rummet. Beteckning $\vb{r} = (x, y, z)$.
